\begin{frame}{Diszjunkt halmaz adatstruktúra}
  \begin{block}{Definíció:}
    A diszjunkt halmaz adatstuktúra diszjunkt halmazok gyűjteményét tartja fenn: $\mathcal{S} = \{ S_1, S_2, \ldots, S_k \}$.
    Minden halmazt egy \textbf{képviselőjével} azonosítjuk, amely eleme ennek a halmaznak.
  \end{block}
  Műveleteket a diszjunkt halmaz adatstruktúrán:
  \begin{itemize}
    \item \textbf{makeset(v)}: létrehoz egy új, egy elemű halmazt: $ \{ v \} $, amelynek v lesz a képviselője. 
    \item \textbf{find(v)}: visszaadja annak a diszjunkt halmaznak a képviselőjét, amelyben v elem van.
    \item \textbf{link(u, v)}: Létrehoz egy új halmazt, ami u és v halmazának az uniója. 
  \end{itemize}
\end{frame}